\documentclass[10pt]{article}
\usepackage{icmlko}

\icmlmeta{IP 파운데이션 월드모델}{189}{한 번도 안 건드린 숫자가 있었다}{열을 늘리는 처방을 여섯 가지 시도했고 전부 졌다. 그러다 능형의 알파가 정식화를 쓴 날부터 기본값 그대로인 것을 봤다 --- 안쪽에서 고르니 세 형태가 다 오른다}

\begin{document}

\twocolumn[
\icmlheader

\begin{abstract}
노트 188이 ``트리는 무리를 곱하지 않고 조건한다''를 남겼다. 조건을 명시로
넣어 봤다 --- 한 무리의 1단 점수를 $K$ 구간으로 자르고 그 구간 표시자를
다른 무리 열에 곱한다. \textbf{여섯 조합 전부 F10 주효과($+$0.3908)
아래다}(제일 좋은 달력 $K{=}2$ 가 $+$0.3869). 이것으로 열을 늘리는 처방이
\textbf{아홉 번 연속} 졌다 --- 짝 곱 · 무리 안 2차 · 2·3차 · 무리 사이 곱 ·
무리 쌓기 · 조건부 쪼갬. 그래서 물었다: 과적합인가 신호가 없는 것인가.
알파를 훑었다. \textbf{규제를 세게 걸면 곱이 아무 일도 안 한다} ---
$\alpha{=}500$ 에서 56열이 $+$0.3891 이고 19열이 $+$0.3888 로 \emph{같다}.
과적합이 아니라 \textbf{신호가 없다.} 그런데 그 훑기가 다른 것을 보여
줬다. \textbf{알파를 1 에서 100 으로 올리기만 해도 판이 $+$0.369 에서
$+$0.382 로 오른다} --- 능형의 ``\texttt{alpha=1.0}''은 정식화를 쓴 날부터
한 번도 안 바뀐 기본값이었다. 판을 보고 고르면 엿보기이므로 \textbf{학습
구간을 시간으로 다시 7 대 3 으로 갈라} 앞쪽으로 적합하고 뒤쪽으로 골랐다 ---
유보는 한 칸도 안 봤다. \textbf{세 형태가 다 오른다}: F6 능형
$+$0.0109($t{=}4.51$) · F11 합집합 $+$0.0130($t{=}4.06$) · F10 도메인별
$+$0.0125($t{=}1.80$). \textbf{열을 늘리는 데 아홉 번 실패하는 동안 안
건드린 숫자 하나가 있었다.} \texttt{lab/tune.py} 로 만들어 도전자 셋을
등록했다. 마지막으로 아티팩트 하나 더. 등록한 도전자가 원본과 소수 넷째
자리까지 같게 나왔다 --- \texttt{pick} 에 공장 대신 객체를 넘겨 안쪽 점수가
전부 nan 이 되고 격자의 첫 값이 뽑힌 것이다. \textbf{노트 187이 하루 전에
세운 규약이 그것을 잡았다} --- ``판이 소수 넷째 자리까지 안 움직이면 바꾼
것이 모형에 안 닿은 것이다.'' 아홉 번째 아티팩트이고 \textbf{규약이 처음으로
값을 했다.}
\end{abstract}
]

\begin{figure*}[t]
\centering
\includegraphics[width=\textwidth]{figs/knob.pdf}
\caption{왼쪽 --- 열을 늘린 시도 아홉, 전부 아래. 가운데 --- 규제를 걸면
곱이 사라진다. 오른쪽 --- 셋 다 오른다.}
\end{figure*}

\section{조건도 안 된다}

\begin{table}[t]
\centering\tsize
\caption{한 무리의 1단 점수를 $K$ 구간으로 잘라 다른 무리에 곱한다.}
\begin{tabular}{lrrr}
\toprule
& 열 & $K{=}2$ & $K{=}3$ \\
\midrule
공유가 조건 & 47/61 & .3468 & .3745 \\
위키가 조건 & 49/64 & .3794 & .3574 \\
달력이 조건 & 45/58 & \textbf{.3869} & .3719 \\
\midrule
\textbf{주효과(기준)} & 19 & \multicolumn{2}{c}{\textbf{.3908}} \\
트리 & 19 & \multicolumn{2}{c}{.4490} \\
\bottomrule
\end{tabular}
\end{table}

\claimx{\textbf{여섯 조합 전부 기준 아래다.} 제일 좋은 것도 $-$0.004
다.}{\textbf{걸어 둔 반대 가능성이 또 맞았다} --- 네 사이클 연속이다(노트
180 · 181 · 188 · 189). \textbf{예측을 걸 때 반대 가능성을 같이 적는 것이
예측 자체보다 잘 맞는다.}}

\section{과적합이 아니라 신호가 없다}

\begin{table}[t]
\centering\tsize
\caption{능형 알파를 훑는다.}
\begin{tabular}{lrr}
\toprule
$\alpha$ & 19열 주효과 & 56열 무리 안 곱 \\
\midrule
1.0 & .3692 & .3706 \\
20.0 & .3752 & .3737 \\
100.0 & \textbf{.3822} & .3784 \\
500.0 & \textbf{.3888} & \textbf{.3891} \\
\bottomrule
\end{tabular}
\end{table}

\claimx{\textbf{$\alpha{=}500$ 에서 둘이 만난다}(.3888 대 .3891). 곱 서른일곱이
있으나 없으나 같다 --- \textbf{규제가 그것들을 0 으로 보내고 아무것도 안
잃는다.}}{그러니 노트 187 · 188 의 실패는 과적합 이야기가 아니다. 낮은
알파에서는 해치고 높은 알파에서는 사라진다 --- \textbf{도움이 되는 알파가
없다.}}

\section{안 건드린 숫자}

\begin{table}[t]
\centering\tsize
\caption{알파를 안쪽 시간 분할(학습 구간 7 대 3)로 고른다. 유보 미접촉.}
\begin{tabular}{lrrrl}
\toprule
& $\alpha$ & 전 & 후 & \\
\midrule
F6 능형 & 1$\to$100 & .3709 & \textbf{.3818} & $t{=}4.51$ \\
F11 합집합 & 1$\to$100 & .3692 & \textbf{.3822} & $t{=}4.06$ \\
F10 도메인별 & 1$\to$500 & .3914 & \textbf{.4039} & $t{=}1.80$ \\
\bottomrule
\end{tabular}
\end{table}

\claimx{\textbf{셋 다 오르고 둘은 확실히 갈린다.} 그리고 결정에 유보 라벨이
한 칸도 안 들어간다 --- 학습 구간을 시간으로 다시 갈라 앞으로 적합하고
뒤로 골랐다.}{\textbf{시간 접힘인 까닭이 있다} --- 노트 182 가 축과 라벨의
관계가 시간에 시드는 것을 보였으므로 무작위 접힘은 낙관적으로 고른다.}

\claimx{\textbf{열을 늘리는 데 아홉 번 실패하는 동안 안 건드린 숫자 하나가
있었다.} 이 흐름에서 노트 170 의 위키 축 전체 이득이 $+$0.0096 이었는데
숫자 하나가 $+$0.0130 이다.}{\texttt{lab/tune.py} 로 만들어
\texttt{F6\_directpool\_tuned} · \texttt{F10\_pershrink\_tuned} ·
\texttt{F11\_poolunion\_tuned} 를 등록했다. \textbf{챔피언 F18 배깅
($+$0.4490)은 아직 안 조정했다} --- 다음이다.}

\section{아홉 번째 아티팩트}

\claimx{\textbf{등록한 도전자가 원본과 소수 넷째 자리까지 같았다}(.3904 대
.3914 · .3708 대 .3709). \texttt{pick} 이 ``정식화를 만드는 함수''를 받는데
객체를 넘겨서, 안쪽 점수가 전부 nan 이 되고 격자의 첫 값이
뽑혔다.}{\textbf{노트 187 이 하루 전에 세운 규약이 잡았다} --- ``판이 소수
넷째 자리까지 안 움직이면 바꾼 것이 모형에 안 닿은 것이다.'' 그 규약을
쓴 것이 \texttt{AXIS\_MODE} 사고였고, \textbf{하루 만에 값을 했다.}}

\claimx{\textbf{같은 모양이 세 번째다} --- 노트 175 의 명시 축
$+$0.0000, 노트 187 의 \texttt{AXIS\_MODE}, 그리고 여기.}{셋 다 ``안
움직였다''를 ``효과가 없다''로 읽을 뻔했고 셋 다 ``안 닿았다''였다.
\textbf{이 판에서 제일 자주 나오는 거짓 음성이다.}}

\begin{table}[t]
\centering\tsize
\caption{남은 것.}
\begin{tabular}{ll}
\toprule
\textbf{F18} & 챔피언의 초매개변수를 안 조정했다 \\
격자 & 알파 일곱 개 --- 더 촘촘히 \\
다른 손잡이 & 접힘 수 · 부스터 수 · 깊이 \\
\bottomrule
\end{tabular}
\end{table}

첫 줄이 곧바로 할 수 있고 값이 클 수 있다. 챔피언 F18 배깅은 부스터 여덟과
기본 깊이로 쓰고 있고 한 번도 안 조정했다 --- 능형이 숫자 하나로 $+$0.013
을 얻었으니 \textbf{챔피언도 그만큼 남아 있을 수 있다.} 다만 조정하면
챔피언이 더 멀어지므로 이 실험실의 ``해석 가능한 모형으로 따라잡기''는 더
어려워진다. \textbf{그래도 잰다} --- 노트 126 이 세운 대로 판정치는 방법을
전제하면 안 된다.

\end{document}
